\scanpage{106}
% L1--5
Donc on peut supposer $E=\complement V$, $V$ \deleted{\unsure{ét|rét.}} ouvert rétrocompact, \deleted{et par suite on aura $\unsure{E-F=F=\complement}$ \unsure{un ouv.}} Il suffit de prouver \inserted{le fermé} $E-F$ constr. dans $X$, i.e. $X-(E-F)=V\cup F$ rétrocompact dans $X$, ce qui est le cas car $V$ et $F$ le sont évidemment \unsure{\ldots}

% Small sketch in the left margin at L3--5.
\figure{\begin{tikzpicture}[scale=.68]
\draw (.1,0)--(.82,2.45); \draw (.08,.28)--(.21,.24) (.63,2.12)--(.75,2.08);
\node at (.47,1.13) {$F$}; \node[above right] at (.8,2.4) {$\complement V$};
\end{tikzpicture}}

% L6
Le cas ``local'' résulte aussitôt du cas global.

% L7--11
\uline{Corollaire} 1 Supposons que de plus, tte partie $U_i$ soit \uline{quasi compacte} \deleted{\unsure{[alors ce}} \uline{pour la topologie constructible} [donc aussi \deleted{pour la top.} compacte pour la top constructible, dans le cas où \deleted{la top} $X$ satisfait (T$_0$)], \inserted{i.e. satisfasse à la condition : \unsure{théorème 1|th. 1}} Alors

\margin{N.B. les hypothèses faites signifient que tte partie fermée \unsure{irrid.|irréd.} de $X$ admet un \deleted{\unsure{pt}} pt générique (unique si (T$_0$) est satisfait \unsure{\ldots})}
% slanted

% L12--14
(iv) Les $U_i^{\mathrm{cons}}$ sont quasi-compacts, \deleted{et \unsure{tout} \inserted{et $X$} \unsure{est} ouvert quasi-compact \unsure{est une base formée d'ouverts} [donc $X$ admet une base formée d'ouverts]}

% L15--18
(v) Soit $E$ une partie \deleted{loc. constructible de} \inserted{fermée de} $X^{\mathrm{cons}}$ [p. ex. une partie loc. constructible \deleted{\unsure{]}} de $X$] alors $E$ \deleted{compact} quasi-compact \inserted{[resp. cpt]} dans $X^{\mathrm{cons}}$ $\Longleftrightarrow$ $E$ quasi-compact dans $X$ $\Longleftrightarrow$ $E\subset$ réunion finie des $U_i$.

% L19--20
(vi) Pour que $X^{\mathrm{cons}}$ soit loc. cpt [cpt] il f. et s. que $X$ satisfasse (T$_0$) [\unsure{resp.} (T$_0$) et soit qu. cpt].

% L21--26
(iv) Résulte de (i) trivialement. Evident $E^{\mathrm{cons}}$ qu. cpt $\Rightarrow$ $E$ qu. cpt en vertu de (i) $\Rightarrow$ $E\subset$ réunion finie \inserted{$V$} des $U_i$ par définition de qu. cpt, donc $E$ fermé dans $V^{\mathrm{cons}}$ [pour la top. induite par $X^{\mathrm{cons}}$] qui est aussi $\bigcup U_i^{\mathrm{cons}}$, et les $U_i^{\mathrm{cons}}$ étant qu. cpts, il s'ensuit que $E$ l'est, O.K. \quad (vi) En résulte \unsure{hic et nunc|l'assertion|trivialement}
% archive sheet number: 105
